\documentclass[a4paper,12pt]{article}
\usepackage[T2A]{fontenc}
\usepackage[utf8]{inputenc}
\usepackage[russian]{babel}
\usepackage{amsmath,amssymb,mathtools}
\usepackage[sfdefault]{noto}
\usepackage{icomma}
\usepackage[margin=21mm,top=23mm,bottom=22mm,headheight=25pt]{geometry}
\usepackage{microtype}
\usepackage{tikz}
\usetikzlibrary{calc,arrows.meta,positioning,shapes.geometric}
\usepackage{booktabs,tabularx,longtable}
\usepackage{enumitem}
\usepackage{parskip}
\usepackage{setspace}
\usepackage{xcolor}
\usepackage{fancyhdr}
\usepackage{hyperref}
\definecolor{accent}{HTML}{167A7A}
\definecolor{darkaccent}{HTML}{155556}
\definecolor{textgray}{HTML}{303A3B}
\definecolor{softline}{HTML}{D5E5E3}
\hypersetup{colorlinks=true,urlcolor=darkaccent,linkcolor=darkaccent,pdftitle={Чек-лист: дробно-рациональные неравенства},pdfauthor={Лёвин Артём Александрович}}
\setstretch{1.13}
\setlength{\parindent}{0pt}
\setlength{\parskip}{4pt}
\setlist[itemize]{leftmargin=*,itemsep=2pt,topsep=4pt}
\setlist[enumerate]{leftmargin=*,itemsep=3pt,topsep=4pt}
\pagestyle{fancy}
\fancyhf{}
\fancyhead[L]{\small\color{textgray}Дробно-рациональные неравенства}
\fancyhead[R]{\small\color{textgray}23.10.26}
\fancyfoot[C]{\small\color{textgray}\thepage}
\renewcommand{\headrulewidth}{0.35pt}
\renewcommand{\headrule}{\hbox to\headwidth{\color{softline}\leaders\hrule height \headrulewidth\hfill}}
\newcommand{\sect}[1]{\vspace{5pt}{\large\bfseries\color{darkaccent}#1}\par\vspace{1pt}}
\newcommand{\key}[1]{\textbf{\color{darkaccent}#1}}
\newcommand{\cb}{\(\square\)}
\newcommand{\RR}{\mathbb{R}}
\newcommand{\sol}{\textbf{Решение.} }
\newcommand{\ans}{\textbf{Ответ:} }
\begin{document}
\thispagestyle{empty}
\begin{center}
\vspace*{32mm}
{\large\color{accent}МАТЕМАТИКА \textbullet\ ПОДГОТОВКА К ОГЭ}\par
\vspace{20mm}
{\Huge\bfseries\color{darkaccent}Чек-лист}\par
\vspace{9mm}
{\LARGE\bfseries Дробно-рациональные\\[2mm] неравенства}\par
\vspace{15mm}
{\large Метод интервалов, ограничения,\\[2mm] сложные числители и замена переменной}\par
\vspace{22mm}
{\large Ярослав Верещагин}\par
\vspace{6mm}
{\normalsize 23 октября 2026 года}\par
\vfill
{\small\color{textgray}Лёвин Артём Александрович\\[2mm]
эксклюзивно для Ярослава Верещагина}\par
\vspace{13mm}
\begin{tikzpicture}[x=1cm,y=1cm]
\draw[accent,line width=1.1pt] (-6.6,0.1) .. controls (-4.2,0.95) and (-2.2,-0.7) .. (0.1,0.06)
  .. controls (2.1,0.7) and (4.2,-0.53) .. (6.6,0.19);
\draw[softline,line width=0.5pt] (-5.8,-0.19) .. controls (-2.3,0.3) and (2.4,-0.34) .. (5.8,-0.09);
\end{tikzpicture}
\end{center}
\clearpage

\sect{Что должно получаться после повторения}
\cb\ Приводить дробно-рациональное неравенство к виду \(\dfrac{P(x)}{Q(x)}\mathrel{\diamond}0\), где \(\diamond\) --- один из знаков \(>,<,\ge,\le\).\\
\cb\ Находить нули числителя и знаменателя и правильно отмечать их на оси.\\
\cb\ Определять знаки дроби на промежутках, проверять граничные случаи.\\
\cb\ Решать задачи с квадратными трёхчленами, постоянным числителем, отсутствием действительных корней и заменой \(t=x^2\).

\sect{1. Основная форма и область допустимых значений}
Любое неравенство с дробью удобно привести к форме
\[
\boxed{\frac{P(x)}{Q(x)}\mathrel{\diamond}0,\qquad Q(x)\ne0.}
\]
Здесь \(P(x)\) --- числитель, \(Q(x)\) --- знаменатель. \key{Сначала перенесите все слагаемые в одну сторону и приведите дроби к общему знаменателю.}

\textbf{Почему нельзя отбросить знаменатель?} Его знак влияет на знак дроби. Умножение неравенства на выражение неизвестного знака без рассмотрения случаев способно изменить ответ.

\textbf{Различайте две группы точек:}
\begin{itemize}
\item \(P(x)=0\) --- дробь обращается в нуль, \emph{если} \(Q(x)\ne0\). Такие точки входят в ответ при знаках \(\ge\) и \(\le\), если подходят по условию.
\item \(Q(x)=0\) --- дробь не определена. Такие точки \key{всегда исключаются}, независимо от строгости неравенства.
\end{itemize}
\textbf{Запись промежутков:} \([a,b]\) включает концы, \((a,b)\) исключает. С бесконечностью всегда круглая скобка: \((\!-\infty,a)\) и \((b,+\infty)\). Несколько промежутков соединяются знаком \(\cup\).

\sect{2. Первый пример: когда требуется отрицательная дробь}
\[
\frac{x-3}{x-5}\le0.
\]
Точки: \(3\) --- нуль числителя, \(5\) --- нуль знаменателя. При \(x=4\) дробь равна \(-1\); это искомый промежуток. Концы: \(3\) включаем, \(5\) исключаем.
\ans\[\boxed{[3,5)}\]

\clearpage
\sect{3. Полный разбор: закрашенные и выколотые точки}
Решим
\[
\frac{x+1}{x-1}\ge0.
\]
\sol Нуль числителя: \(x=-1\); нуль знаменателя: \(x=1\), значит \(x\ne1\). Наносим точки на ось. Пробное значение \(x=0\) находится между ними:
\[
\left.\frac{x+1}{x-1}\right|_{x=0}=\frac{1}{-1}=-1<0.
\]
По обе стороны от простых корней знак меняется; получаем следующую схему.
\begin{center}
\begin{tikzpicture}[x=0.82cm,y=0.7cm,>=Latex]
 \draw[->,line width=.9pt] (0,0)--(13,0) node[right] {$x$};
 \draw[very thick,accent] (1.2,0.12)--(4.1,0.12);
 \draw[very thick,accent] (9.0,0.12)--(12.4,0.12);
 \draw[fill=accent,draw=accent] (4.1,0) circle (3.4pt);
 \draw[fill=white,draw=darkaccent,line width=1pt] (9.0,0) circle (3.4pt);
 \node[below=5pt] at (4.1,0) {$-1$};
 \node[below=5pt] at (9,0) {$1$};
 \node[above=5pt] at (2.1,.12) {$+$};
 \node[above=5pt] at (6.6,.12) {$-$};
 \node[above=5pt] at (10.9,.12) {$+$};
\end{tikzpicture}
\end{center}
\ans \(\boxed{(-\infty,-1]\cup(1,+\infty)}\).

\textbf{Самопроверка:} число \(-1\) включается, потому что дробь равна нулю и задан знак \(\ge\); число \(1\) исключается, поскольку знаменатель равен нулю.

\sect{4. Квадратные трёхчлены: сначала найдите корни}
Рассмотрим пример с четырьмя граничными точками:
\[
\frac{x^2-5x+6}{x^2-12x+35}\ge0.
\]
Разложим числитель и знаменатель на множители:
\[
\frac{(x-2)(x-3)}{(x-5)(x-7)}\ge0,\qquad x\ne5,7.
\]
Нули числителя \(2\) и \(3\); нули знаменателя \(5\) и \(7\). Все четыре корня простые. Между ними получаем знаки, проверяя, например, \(x=0\):
\[
\left.\frac{(x-2)(x-3)}{(x-5)(x-7)}\right|_{x=0}
=\frac{6}{35}>0.
\]
\begin{center}
\renewcommand{\arraystretch}{1.23}
\begin{tabular}{@{}lccccc@{}}
\toprule
Промежуток & \((\!-\infty,2)\) & \((2,3)\) & \((3,5)\) & \((5,7)\) & \((7,+\infty)\)\\
\midrule
Знак дроби & \(+\) & \(-\) & \(+\) & \(-\) & \(+\)\\
\bottomrule
\end{tabular}
\end{center}
\ans \(\boxed{(-\infty,2]\cup[3,5)\cup(7,+\infty)}\).

\textbf{Важно:} знак при пересечении \emph{простого} корня множителя меняется. Если кратность корня чётная, знак произведения при переходе через такую точку \emph{сохраняется}. Наиболее надёжный подход --- определить знак на каждом промежутке пробным значением или учитывать кратности всех множителей.

\clearpage
\sect{5. Если справа стоит число: сначала привести к нулю}
Разберём
\[
\frac{2x^2+18x-4}{x^2+9x+8}>2.
\]
\textbf{Шаг 1.} Запишите ограничения: \(x^2+9x+8\ne0\).\\
\textbf{Шаг 2.} Перенесите \(2\) влево и приведите к одному знаменателю:
\begin{align*}
\frac{2x^2+18x-4}{x^2+9x+8}-2&>0,\\
\frac{2x^2+18x-4-2(x^2+9x+8)}{x^2+9x+8}&>0,\\
\frac{-20}{x^2+9x+8}&>0,\\
\frac{-20}{(x+1)(x+8)}&>0.
\end{align*}
\textbf{Шаг 3.} В числителе постоянное отрицательное число; нулей числителя нет. Знаменатель обращается в нуль при \(x=-8\) и \(x=-1\). Дробь положительна там, где знаменатель отрицателен:
\[
(x+1)(x+8)<0\quad\Longrightarrow\quad-8<x<-1.
\]
\ans \(\boxed{(-8,-1)}\).

\textbf{Ключевая проверка:} пробное значение подставляют уже в \emph{преобразованную} дробь. Например, при \(x=-4\) имеем \(-20/(-12)>0\). Точки \(-8\) и \(-1\) в ответ не входят.

\sect{6. Особые случаи, которые не нарушают алгоритм}
\textbf{А. У числителя нет действительных корней.} Например,
\[
x^2+x+1=0,\quad D=1-4=-3<0.
\]
Так как старший коэффициент положителен, \(x^2+x+1>0\) при всех \(x\in\RR\). На ось нули этого многочлена наносить не нужно.

\textbf{Б. Числитель --- ненулевая константа.} Например, \(-20\) никогда не обращается в нуль. Выясняйте знак знаменателя.

\textbf{В. Пробное число совпадает с корнем.} Нельзя подставлять \(x=0\), если \(0\) --- нуль числителя или знаменателя. Возьмите другое число \emph{внутри исследуемого промежутка}, например \(x=1\).

\textbf{Г. Один и тот же множитель сверху и снизу.} Даже после сокращения нужно сохранить запрет на исходный нуль знаменателя:
\[
\frac{(x-3)(x+2)}{(x-3)(x-1)}\ge0,
\qquad x\ne1,\;x\ne3.
\]
При допустимых \(x\) дробь равна \((x+2)/(x-1)\), но число \(3\) по-прежнему выколото. Поэтому
\ans\[\boxed{(-\infty,-2]\cup(1,3)\cup(3,+\infty)}\]
\textbf{Д. Корень чётной кратности.} У множителя \((x-a)^2\) знак не меняется при переходе через \(a\). Точку \(a\) всё равно проверяют на включение, если она допустима и знак неравенства нестрогий.

\sect{7. Сложный числитель и замена \(t=x^2\)}
Сначала ищем нули: \(P(x)=0\), \(Q(x)=0\). При этом уравнение \(AB+1=0\) \key{нельзя} заменять системой \(A=0\) или \(B=0\). Правило нулевого произведения работает только для \(AB=0\).

Для выражения \(x^4-2x^2-8\) удобно положить \(t=x^2\), где \(t\ge0\):
\[
t^2-2t-8=0 \iff (t-4)(t+2)=0.
\]
Из \(t=4\) имеем \(x=\pm2\); значение \(t=-2\) отвергаем, так как \(x^2\ge0\).

Другой характерный пример:
\[
x^4+x^2+1=0
\iff t^2+t+1=0,\quad D=-3<0.
\]
Действительных корней здесь нет. Более того, \(x^4+x^2+1>0\) для всякого действительного \(x\).

\sect{8. Разбор многочлена третьей степени}
Решим
\[
\frac{x^3+x^2+x}{9x^2-25}\ge0.
\]
Числитель раскладывается так:
\[
x^3+x^2+x=x(x^2+x+1).
\]
Так как \(x^2+x+1>0\), знак числителя определяется знаком \(x\). В знаменателе
\[
9x^2-25=(3x-5)(3x+5).
\]
Критические точки: \(-\frac53\) и \(\frac53\) из знаменателя (выколоты), \(0\) из числителя (включается). Калибровку можно сделать при \(x=1\):
\[
\frac{1+1+1}{9-25}=-\frac{3}{16}<0.
\]
Знаки на промежутках слева направо: \(-,+,-,+\). Поэтому
\ans\[\boxed{\left(-\frac53,0\right]\cup\left(\frac53,+\infty\right)}\]

\sect{9. Контрольный список перед записью ответа}
\cb\ Справа \(0\); дробь приведена к единому знаменателю.\\
\cb\ Учтено ограничение \(Q(x)\ne0\) из \emph{исходной} записи.\\
\cb\ Все действительные корни числителя и знаменателя найдены.\\
\cb\ Точки отсортированы на оси, знаки проверены с учётом кратности.\\
\cb\ Знак \(\ge/\le\) допускает нули дроби, но никогда нули знаменателя.\\
\cb\ Промежутки объединены через \(\cup\); скобки и бесконечности оформлены верно.

\sect{10. Почему знак не всегда чередуется}
Сравните метод интервалов с задачей, в которой один из нулей имеет чётную кратность:
\[
\frac{(x+1)^2(x-4)}{x-2}>0.
\]
Знаменатель обращается в нуль при \(x=2\); числитель --- при \(x=-1\) и \(x=4\). Точка \(-1\) является \emph{двойным нулём}.

При \(x=0\) значение дроби равно \(\frac{1\cdot(-4)}{-2}=2>0\). Поэтому знаки на интервалах располагаются так:
\begin{center}
\renewcommand{\arraystretch}{1.2}
\begin{tabular}{@{}lcccc@{}}
\toprule
Интервал & \((\!-\infty,-1)\) & \((-1,2)\) & \((2,4)\) & \((4,+\infty)\)\\
\midrule
Знак & \(+\) & \(+\) & \(-\) & \(+\)\\
\bottomrule
\end{tabular}
\end{center}
Переход через \(-1\) сохраняет знак. При этом \(-1\) не входит в ответ: неравенство \emph{строгое}, а дробь в этой точке равна нулю. Через \(2\) и \(4\) знак меняется, поскольку это простые корни соответствующих множителей.
\[
\boxed{(-\infty,-1)\cup(-1,2)\cup(4,+\infty)}.
\]
\textbf{Практическое правило:} если сомневаетесь в чередовании, вычислите знак на каждом промежутке отдельно. Это избавляет от ошибок с повторными корнями.

\clearpage
\thispagestyle{fancy}
\begin{center}
{\Large\bfseries\color{darkaccent}Алгоритм: дробно-рациональное неравенство}
\end{center}
\vspace{5mm}
\begin{center}
\begin{tikzpicture}[>=Latex, node distance=8mm, font=\small,
 startstop/.style={ellipse,draw=darkaccent,fill=softline!35,align=center,minimum width=40mm,minimum height=10mm},
 stage/.style={rectangle,rounded corners=2pt,draw=accent,align=center,text width=100mm,minimum height=12mm,inner sep=4pt},
 decision/.style={diamond,aspect=2.4,draw=darkaccent,align=center,text width=37mm,inner sep=2pt},
 aside/.style={rectangle,rounded corners=2pt,draw=accent,align=center,text width=45mm,minimum height=13mm,inner sep=3pt},
 every path/.style={line width=.9pt}]
\node[startstop] (start) {Начало};
\node[stage,below=of start] (normal) {Перенесите всё влево: $\dfrac{P(x)}{Q(x)}\mathrel{\diamond}0$.\\ Сохраните ОДЗ: $Q(x)\ne0$};
\node[stage,below=of normal] (factor) {Разложите $P(x)$ и $Q(x)$ на множители.\\ Найдите их действительные нули};
\node[decision,below=10mm of factor] (root) {Есть нули\\ числителя?};
\node[aside,at={($(root.south)+(-39mm,-25mm)$)}] (yes) {Нули числителя:\\ отметьте на оси;\\ проверьте включение};
\node[aside,at={($(root.south)+(39mm,-25mm)$)}] (no) {Нулей числителя нет};
\node[stage,below=47mm of root] (denom) {Нули знаменателя отметьте выколотыми;\\ разбейте ось на промежутки};
\node[stage,below=of denom] (signs) {Определите знак дроби на промежутках.\\ Учитывайте кратности; при сомнении проверьте пробным числом};
\node[stage,below=of signs] (select) {Выберите промежутки нужного знака;\\ проверьте допустимые нули числителя};
\node[stage,below=of select] (write) {Запишите ответ объединением промежутков\\ с правильными скобками};
\node[startstop,below=of write] (finish) {Конец};
\draw[->] (start)--(normal);
\draw[->] (normal)--(factor);
\draw[->] (factor)--(root);
\draw[->] (root.west)-|node[pos=.25,above]{да}(yes.north);
\draw[->] (root.east)-|node[pos=.25,above]{нет}(no.north);
\coordinate (join) at ($(denom.north)+(0,11mm)$);
\draw[-] (yes.south) |- (join);
\draw[-] (no.south) |- (join);
\draw[->] (join)--(denom.north);
\draw[->] (denom)--(signs);
\draw[->] (signs)--(select);
\draw[->] (select)--(write);
\draw[->] (write)--(finish);
\end{tikzpicture}
\end{center}
\clearpage

\sect{Тренировочный блок --- Путь А}
\textbf{10 заданий.} Решите неравенства методом интервалов. Постарайтесь в каждом задании отдельно указать нули числителя, запрещённые значения и нужные промежутки. Сложность постепенно возрастает.
\begin{enumerate}[label=\textbf{\arabic*.},leftmargin=10mm,itemsep=11pt]
\item \(\displaystyle\frac{x+2}{x-4}\ge0.\)
\item \(\displaystyle\frac{x-5}{x+1}\le0.\)
\item \(\displaystyle\frac{x^2-5x+4}{x^2-9x+18}\ge0.\)
\item \(\displaystyle\frac{2x^2+8x-3}{x^2+4x+3}>2.\)
\item \(\displaystyle\frac{x^2+4}{x^2-9}\le0.\)
\item \(\displaystyle\frac{-6}{x^2-5x+6}\ge0.\)
\item \(\displaystyle\frac{x^3+x^2+x}{4x^2-9}\ge0.\)
\item \(\displaystyle\frac{(x-1)^2(x+2)}{x-3}<0.\)
\item \(\displaystyle\frac{(x-3)(x+2)}{(x-3)(x-1)}\ge0.\)
\item \(\displaystyle\frac{x^4-2x^2-8}{x^2+1}>0.\)
\end{enumerate}
\vfill
\textbf{Сосредоточьтесь на случаях 4, 8, 9, 10.} Здесь соответственно проверяются приведение к нулю, корень чётной кратности, запрет после сокращения и замена \(t=x^2\).
\clearpage

\sect{Ответы к тренировочному блоку}
\textbf{Проверьте сначала свои скобки и исключённые точки, затем сравнивайте интервалы.}
\vspace{5mm}
\renewcommand{\arraystretch}{1.52}
\begin{longtable}{@{}p{11mm}p{82mm}p{54mm}@{}}
\toprule
\textbf{№} & \textbf{Ответ} & \textbf{Ориентир для проверки}\\
\midrule
\endfirsthead
\toprule
\textbf{№} & \textbf{Ответ} & \textbf{Ориентир для проверки}\\
\midrule
\endhead
1 & \( (-\infty,-2]\cup(4,+\infty) \) & \(x=4\) исключить\\
2 & \( (-1,5] \) & \(x=-1\) исключить\\
3 & \( (-\infty,1]\cup(3,4]\cup(6,+\infty) \) & Нули: \(1,4\); запрет: \(3,6\)\\
4 & \( (-3,-1) \) & Числитель: \(-9\)\\
5 & \( (-3,3) \) & Числитель \(>0\)\\
6 & \( (2,3) \) & Нуль числителя отсутствует\\
7 & \( (-\frac32,0]\cup(\frac32,+\infty) \) & \(x^2+x+1>0\)\\
8 & \( (-2,1)\cup(1,3) \) & \(x=1\) --- двойной нуль\\
9 & \( (-\infty,-2]\cup(1,3)\cup(3,+\infty) \) & \(x=3\) остаётся запрещённым\\
10 & \( (-\infty,-2)\cup(2,+\infty) \) & \(t=4\), а \(t=-2\) недопустимо\\
\bottomrule
\end{longtable}
\vfill
\begin{center}
{\small\color{textgray}Лёвин Артём Александрович \textbullet\ эксклюзивно для Ярослава Верещагина}
\end{center}
\end{document}
